\section{Why point ideals do not give the required eightfold character}
\label{sec:spin-obstructions}
This section excludes every external product of point-ideal characters in dimension eight as a source of a nonzero Weil component on the full family. It checks the whole Chern character under every Fourier--Mukai transform. The obstruction is a property of their rational Clifford orbits, rather than a failure of a particular transform.

If $R=H^1(A,\Q)$, the cohomological spin representation of the split quadratic space $R\oplus R^\vee$ is $\bigwedge^*R$. The Lie algebra acts through
\begin{equation}\label{eq:spin-lie-action}
 \mathfrak{so}(R\oplus R^\vee)
 =\bigwedge^2R\oplus\mathfrak{gl}(R)\oplus\bigwedge^2R^\vee.
\end{equation}
The three actions are exterior multiplication, the derivation induced by an endomorphism minus half its trace, and contraction. The Clifford compatibility of Fourier--Mukai equivalences is established in Golyshev--Lunts--Orlov \cite[\S3 and Proposition~4.3.7, author pp.~17--19]{GLO}; the relative construction in its proof applies between distinct abelian varieties. Consequently such an equivalence transports the Lie stabilizer of a character by a rational isomorphism. Multiplication by $e^B$ for a rational two-form is a Clifford shear and does the same. Isogeny pullback is the exterior map induced by the rational isomorphism on $H^1$, and isogeny pushforward is a nonzero scalar times its inverse. These operations also preserve the isomorphism type of the stabilizer.

\begin{proposition}[Products of two point ideals]\label{prop:point-product-spin}
Let $X,Y$ be abelian fourfolds, let $p\in X$, $q\in Y$, and let $L,M$ be line bundles. No character obtained from
\[
 (I_p\otimes L)\boxtimes(I_q\otimes M)
\]
by Fourier--Mukai equivalences, rational isogeny pullbacks or pushforwards, rational $B$-field transformations and nonzero scalar multiplication remains Hodge on a full polarized Weil domain of signature $(4,4)$.
\end{proposition}

\begin{proof}
Tensoring by $L$ and $M$ can be removed by an exponential transformation. Put $U=H^1(X,\Q)$, $V=H^1(Y,\Q)$, and let $u\in\bigwedge^8U$, $v\in\bigwedge^8V$ be the point classes. The character is
\[
 f=(1-u)(1-v)=1-u-v+uv.
\]
Its Lie stabilizer in \eqref{eq:spin-lie-action} is exactly
\begin{equation}\label{eq:point-product-stabilizer}
 \mathfrak{sl}(U)\oplus\mathfrak{sl}(V).
\end{equation}
To check this, let $(B_2,D,\gamma)$ annihilate $f$. Its degree-two part is $B_2$, so $B_2=0$. Its degree-fourteen part is $\iota_\gamma(uv)$; contraction with a top volume identifies $\bigwedge^2(U\oplus V)^\vee$ with $\bigwedge^{14}(U\oplus V)$, so $\gamma=0$. The degree-zero term then forces $\operatorname{tr}D=0$. The two off-diagonal blocks of $D$ acting on $u+v$ have distinct bidegrees $(7,1)$ and $(1,7)$. Each map is injective by contraction with the corresponding volume, so these blocks vanish. The remaining degree-eight term forces the two diagonal traces separately to vanish. Conversely these conditions fix every term of $f$, proving \eqref{eq:point-product-stabilizer}.

Suppose a transformed character $z$ were Hodge throughout a full Weil domain. The Hodge circles at its points are the conjugates of a noncentral circle in $SU(4,4)$. Their connected normal closure is all of $SU(4,4)$, because its real Lie algebra is simple. Thus this Lie algebra lies in the spin stabilizer of $z$. By the preceding rational conjugacy it would embed into
\[
 \mathfrak{sl}_8(\R)\oplus\mathfrak{sl}_8(\R).
\]
Each projection is either zero or injective, by simplicity. A nonzero projection would be an isomorphism, since both source and target have real dimension $63$. But their real ranks are $4$ and $7$, respectively. Every projection must therefore be zero, a contradiction. The same reasoning applies if the character is Hodge on a nonempty analytic open subset, by analytic continuation of the flat Hodge condition.
\end{proof}

The four nonzero coefficients in the expansion of $f$ were the only property of their values used in the stabilizer calculation. Thus the conclusion also holds for products with factor characters $r-au$ and $s-bv$, when $r,a,s,b$ are all nonzero.

\begin{proposition}[A single point secant]\label{prop:single-point-spin}
Let $Z$ be an abelian eightfold and let $c\in\Q^\times$. Apply any sequence of the rational Clifford operations above to $1+c[\mathrm{pt}_Z]$, obtaining $z$. If $z$ remains Hodge on a full polarized Weil domain of signature $(4,4)$ with polarization $\theta$, then
\[
 z\in\Q[\theta].
\]
In particular its Weil component is zero.
\end{proposition}

\begin{proof}
Let $R=H^1(Z,\Q)$, of dimension sixteen. The degree-two and degree-fourteen terms of \eqref{eq:spin-lie-action} force the exterior and contraction blocks of an element stabilizing $1+c[\mathrm{pt}]$ to vanish. Its degree-zero part forces its middle block to have trace zero. Hence the stabilizer is exactly $\mathfrak{sl}(R)$, and it fixes the two rational pure spinors $1$ and $[\mathrm{pt}]$ separately.

A nonzero spinor is called pure when its Clifford annihilator is a maximal isotropic subspace. The rational operations in the statement preserve this property. After transporting the two endpoint spinors, the connected stabilizer of $z$ therefore fixes two rational pure spinors separately, and $z$ is their linear combination. The full Weil Hodge condition puts $SU(h)$ in that connected stabilizer, by the Hodge-circle argument above.

We classify its invariant rational pure spinors. Write $H=H^1(A,\Q)$ for the target. This rational representation of $SU(h)$ is irreducible. Over the imaginary quadratic field its two absolutely irreducible constituents are the standard and dual eight-dimensional representations, which are nonisomorphic and exchanged by descent. Thus no proper rational subrepresentation can be formed from them.

If $p$ is an invariant rational pure spinor, its annihilator $L_p\subset H\oplus H^\vee$ is a rational invariant maximal isotropic subspace of dimension sixteen. The image under projection to $H^\vee$ is either zero or all of $H^\vee$, by irreducibility. In the first case $L_p=H$, and $p$ is a scalar multiple of the top spinor $[\mathrm{pt}_A]$. In the second case the projection is an isomorphism, so $L_p$ is the graph of a map $H^\vee\to H$. Isotropy makes this map an alternating two-form $D\in\bigwedge^2H$, and invariance makes it $SU(h)$-invariant.

The invariant part of $\bigwedge^2H$ is $\Q\theta$. Indeed after complexification $H=V\oplus V^*$, and the only invariant in
\[
 \bigwedge^2V\oplus(V\otimes V^*)\oplus\bigwedge^2V^*
\]
is the scalar pairing in the middle term. Therefore $D=a\theta$ for a rational $a$. The pure-spinor line annihilated by its graph is that of $e^{a\theta}$, after adjusting the graph sign convention. Every invariant rational pure spinor consequently lies in $\Q[\theta]$, so the same is true of $z$.
\end{proof}

For the particular character $1-[\mathrm{pt}]$ of a point ideal, a pure nonzero Weil target is already excluded by the Euler pairing. Its self-pairing is $-2$, whereas $\int_Aw^2>0$ for a nonzero real middle Weil class on an eightfold. Derived equivalences and $B$-fields preserve this pairing; isogenies scale it by positive factors. Proposition~\ref{prop:single-point-spin} is stronger because it also permits an arbitrary polynomial in the polarization alongside the proposed Weil component.


\begin{theorem}[All products of point secants]\label{thm:all-point-products}
Let $A_1,\ldots,A_s$ be positive-dimensional abelian varieties with
$\sum_i\dim A_i=8$. For nonzero rational $r_i,c_i$, put
\[
 f=\boxtimes_{i=1}^s\bigl(r_i+c_i[\mathrm{pt}_{A_i}]\bigr).
\]
Apply any sequence of Fourier--Mukai equivalences, rational isogeny
pullbacks or pushforwards, rational $B$-field transformations and
nonzero scalar multiplication. If the resulting character $z$ is
Hodge on a full polarized Weil domain of signature $(4,4)$, then
$z\in\Q[\theta]$. Thus its Weil component is zero.
In particular no external product of line-bundle twists of point
ideals, for any partition of dimension eight, gives the required
full-family Weil detector character by these operations.
\end{theorem}

\begin{proof}
Rescale each factor to write its spinor as
$f_i=1+b_i\operatorname{vol}_i$, with $b_i\ne0$.
Let $g_i=\dim A_i$, $R_i=H^1(A_i,\R)$ and
$T_i=R_i\oplus R_i^\vee$. Then $\dim_\R T_i=4g_i$.
Under the rational Clifford operations, the action of $SU(4,4)$ on
the total generalized vector space is conjugate to its usual action
on $H^1(A,\R)\oplus H^1(A,\R)^\vee$. The standard real
representation $W=H^1(A,\R)$ is irreducible of dimension sixteen;
its dual is isomorphic to it through the polarization. Hence the
total representation is $W\oplus W$, and every invariant real
subspace has dimension a multiple of sixteen.

First suppose that $k>0$ of the factors are elliptic. Their spinors
are pure and have Clifford annihilators of real dimension two.
For $g_i\ge2$, Clifford multiplication is the map
\[
 R_i\oplus R_i^\vee\longrightarrow\bigwedge^{\mathrm{odd}}R_i,
 \qquad(x,\xi)\longmapsto x+b_i\iota_\xi\operatorname{vol}_i.
\]
Its two summands have degrees one and $2g_i-1$, which are distinct,
so it is injective. In the external product, the images from the
different factors lie in distinct parity sectors. Its annihilator
is therefore the direct sum of the factor annihilators and has
dimension $2k$. If $z$ is Hodge on the full domain, this annihilator
is invariant under the conjugate of $SU(4,4)$, so $2k$ must be a
multiple of sixteen. Since $0<k\le8$, this forces $k=8$.
All factors are then elliptic and $f$, and hence $z$, is pure.
The invariant-pure-spinor argument in the proof of
Proposition~\ref{prop:single-point-spin} gives $z\in\Q[\theta]$.

It remains to treat $g_i\ge2$ for every factor. The orthogonal Lie
algebra has block decomposition
\[
 \mathfrak{so}\Bigl(\bigoplus_iT_i\Bigr)
 =\bigoplus_i\mathfrak{so}(T_i)
  \ \oplus\!\bigoplus_{i<j}(T_i\otimes T_j).
\]
A mixed block $T_i\otimes T_j$ acts on $f$ through the tensor product
of the two injective Clifford maps just displayed. It is therefore
injective. Its image has odd parity in precisely factors $i,j$;
these parity sectors are distinct for different pairs and are
disjoint from the images of the diagonal Lie blocks. Consequently
every Lie stabilizer element has zero mixed blocks. The full
$SU(4,4)$ action therefore preserves each $T_i$ separately.

Each $4g_i$ must now be a multiple of sixteen. A partition of eight
into positive multiples of four is either $8$ or $4+4$.
The first case is Proposition~\ref{prop:single-point-spin}; the
second is Proposition~\ref{prop:point-product-spin} and its
nonzero-coefficient extension. They give the stated conclusion.
Line-bundle twists can be removed by exponential transformations,
which were included among the permitted operations.
\end{proof}

The proof deliberately does not identify the internal stabilizer
of a surface factor with a special linear group. In dimension two,
exterior multiplication in degree two can cancel contraction of the
top term. The injectivity of the odd Clifford map is the property
needed for the mixed-block argument, and it holds in every factor
of dimension at least two.

These exclusions leave the pure character of Proposition~\ref{prop:sixteen-line-character} untouched. Its two complex determinant lines are exchanged by the imaginary quadratic field and are not rational pure-spinor endpoints. They also leave arbitrary full-support sheaves and other perfect complexes open. A positive construction must meet both requirements: a whole character that stays Hodge on the full domain and an invariant obstruction space satisfying Theorem~\ref{thm:pure-weil-ext} when the normalized character is pure.
